\documentclass[12pt]{amsart}
%prepared in AMSLaTeX, under LaTeX2e
\addtolength{\oddsidemargin}{-.65in} 
\addtolength{\evensidemargin}{-.65in}
\addtolength{\topmargin}{-.55in}
\addtolength{\textwidth}{1.3in}
\addtolength{\textheight}{.9in}

\renewcommand{\baselinestretch}{1.05}

\usepackage{verbatim} % for "comment" environment

\usepackage{palatino}

\newtheorem*{thm}{Theorem}
\newtheorem*{defn}{Definition}
\newtheorem*{example}{Example}
\newtheorem*{problem}{Problem}
\newtheorem*{remark}{Remark}

\usepackage{fancyvrb,xspace,dsfont}

\usepackage[final]{graphicx}

% macros
\usepackage{amssymb}
\usepackage[dvipsnames]{xcolor}

\usepackage{hyperref}
\hypersetup{pdfauthor={Ed Bueler},
            pdfcreator={pdflatex},
            colorlinks=true,
            citecolor=blue,
            linkcolor=red,
            urlcolor=blue,
            }

\newcommand{\bn}{\mathbf{n}}
\newcommand{\br}{\mathbf{r}}
\newcommand{\bv}{\mathbf{v}}
\newcommand{\bx}{\mathbf{x}}
\newcommand{\by}{\mathbf{y}}

\newcommand{\bJ}{\mathbf{J}}
\newcommand{\bV}{\mathbf{V}}

\newcommand{\cB}{\mathcal{B}}
\newcommand{\cD}{\mathcal{D}}
\newcommand{\cE}{\mathcal{E}}
\newcommand{\cF}{\mathcal{F}}
\newcommand{\cH}{\mathcal{H}}
\newcommand{\cL}{\mathcal{L}}
\newcommand{\cM}{\mathcal{M}}
\newcommand{\cR}{\mathcal{R}}
\newcommand{\cV}{\mathcal{V}}
\newcommand{\cW}{\mathcal{W}}

\newcommand{\CC}{\mathbb{C}}
\newcommand{\NN}{\mathbb{N}}
\newcommand{\RR}{\mathbb{R}}
\newcommand{\ZZ}{\mathbb{Z}}

\renewcommand{\Im}{\mathrm{Im}}
\renewcommand{\Re}{\mathrm{Re}}

\newcommand{\eps}{\epsilon}
\newcommand{\lam}{\lambda}
\newcommand{\lap}{\triangle}

\newcommand{\Div}{\nabla\cdot}
\newcommand{\grad}{\nabla}

\newcommand{\one}{\mathds{1}}

\newcommand{\ip}[2]{\ensuremath{\left<#1,#2\right>}}

\newcommand{\image}{\operatorname{im}}
\newcommand{\onull}{\operatorname{null}}
\newcommand{\rank}{\operatorname{rank}}
\newcommand{\range}{\operatorname{range}}
\newcommand{\trace}{\operatorname{tr}}
\newcommand{\Span}{\operatorname{span}}

\newcommand{\prob}[1]{\bigskip\noindent\textbf{#1.}\quad }

\newcommand{\pts}[1]{(\emph{#1 pts}) }
\newcommand{\epart}[1]{\medskip\noindent\textbf{(#1)}\quad }
\newcommand{\ppart}[1]{\,\textbf{(#1)}\quad }

\newcommand{\ds}{\displaystyle}

\newcommand{\nex}{\medskip\noindent}


\begin{document}
\scriptsize \noindent Math 617 Functional Analysis (Bueler) \hfill \emph{assigned 16 March 2026; {\color{red} version 2}}
\normalsize\medskip

\Large\centerline{\textbf{Assignment 6}}
\large
\medskip

\centerline{\textbf{Due Friday 27 March at the beginning of class}}
\medskip
\normalsize

\thispagestyle{empty}

\bigskip
\noindent This Assignment is based on Chapters 3 and 5 of our textbook,\footnote{K.~Saxe,~\emph{Beginning Functional Analysis}, Springer 2010.} and on the slides, especially those for \href{https://bueler.github.io/fa/assets/slides/S26/week10.pdf}{weeks 10 \& 11}.

\newcommand{\augment}[1]{\quad\,\strut \begin{minipage}[t]{0.7\textwidth} \emph{#1}\end{minipage}}

\newcommand{\augpart}[1]{\emph{\textbf{(#1)}}}

\bigskip
\noindent \textsc{Do the following} Exercises from Chapter 5 (see pages 129--134):

\begin{itemize}
\item \textbf{Exercise 5.1.1}
\item \textbf{Exercise 5.2.1}
\end{itemize}

\medskip
\noindent \textsc{Do the following additional problems.}

\prob{P12}  \emph{Use the classical definition of absolute continuity, and the FTC2, from the slides.}

\medskip \noindent Suppose throughout that $f,g:[a,b] \to \RR$ are absolutely continuous.

\epart{a} Show that $fg$ is also absolutely continuous.

\epart{b} Show that the product rule $\big(f(x) g(x)\big)' = f'(x) g(x) + f(x) g'(x)$ holds almost everywhere.

\epart{c} Show the following integration-by-parts formula, justifying the existence of all (Lebesgue) integrals and the existence of point values:
	$$\int_a^b f(x) g'(x)\,dx = f(b)g(b) - f(a)g(a) - \int_a^b f'(x) g(x)\,dx$$


\prob{P13}  \emph{Use the product rule for partial derivatives of scalar functions, as stated in the slides.}

\epart{a} Assume that $\Omega\subset \RR^d$ is open.  Suppose that $f:\Omega\to\RR$ and $\bV:\Omega\to \RR^d$ are differentiable at every point of $\Omega$.  Using the definition of the divergence and gradient, show this product rule:
    $$\Div (f\,\bV) = \grad f\cdot \bV + f\,\Div \bV.$$

\epart{b} Assume that $\Omega = (0,1)^d \subset \RR^d$ is the open unit cube.  Suppose that $f,g\in C_c^1(\Omega)$.  Show the ``integration-by-parts 1: partial derivatives'' rule in the slides, namely
    $$\int_\Omega \frac{\partial f}{\partial x_k} g\,dm = -\int_\Omega f \frac{\partial g}{\partial x_k}\,dm,$$
with careful attention to why there is no integral over $\partial \Omega$ in this rule.


\prob{P14}  \emph{Basic properties of the Fourier transform make good exercises with the $L^1(\RR^1)$ space.  There is a more advanced analysis, the Plancherel theorem, which shows that in fact the transform is a Hilbert space isometric isomorphism of $L^2(\RR^1)$, but that would take us too far afield.}

\medskip \noindent Suppose $f:\RR^1\to\CC$ is measurable.  The \emph{Fourier transform} of $f$ is
\begin{equation}
\hat f(t) = \frac{1}{\sqrt{2\pi}} \int_{-\infty}^\infty f(x) \, e^{-ixt} \, dx, \label{ft}
\end{equation}
for those $t\in\RR^1$ such that the integral exists, in which case $\hat f(t)\in\CC$.

\epart{a} Suppose $f\in L^1(\RR^1)$.  Show that $\hat f(t)$ is well-defined for every $t$, and that $\|\hat f\|_\infty \le \frac{1}{\sqrt{2\pi}} \|f\|_1$.  (\emph{Thus} $\hat f\in L^\infty(\RR^1)$.)

\epart{b} Again suppose $f\in L^1(\RR^1)$.  Use Lebesgue's dominated convergence theorem (DCT) to show that $\hat f$ is uniformly continuous.  (\emph{Thus} $\hat f\in C(\RR^1)\cap L^\infty(\RR^1)$.)

\epart{c} Suppose $f\in L^1(\RR^1)$, and also that $g(x)=-ixf(x)$ is in $L^1(\RR^1)$.  Use the DCT to show that $\hat f$ is differentiable at every $t$, and that $\hat f'(t) = \hat g(t)$.  (\emph{Thus} $\hat f\in C^1(\RR^1)\cap L^\infty(\RR^1)$.)


\prob{P15}  Suppose $f\in L^1([0,1])$ and define the \emph{Ces\`aro operator}
	$$(Tf)(s) = \frac{1}{s} \int_0^s f(t)\,dt.$$

\epart{a}  Show that $T$ is linear, and determine a kernel $k(s,t)$ so that $T$ can be written as a Fredholm integral operator, in the form $(Tf)(s) = \int_0^1 k(s,t)\,f(t)\,dt$.

\epart{b}  Show that $(Tf)(s)$ is continuous at $s>0$.  (\emph{Hint.}  DCT.)

\epart{c}  Give an example of $f\in L^1([0,1])$ which shows that $Tf$ need not be continuous at $s=0$.  The same example should show that $Tf\notin L^\infty([0,1])$.


\prob{P16}  Suppose $\rho=\rho(t,x)=\rho(t,x_1,x_2,x_3)$, for $t\in\RR^1$ and $x\in\RR^3$, is continuously-differentiable (in each variable).  Suppose $\bJ=\bJ(t,x) = \bJ(t,x_1,x_2,x_3)$ is also continuously-differentiable.  We may interpret $\rho$ as a scalar mass density, or charge density etc., and $\bJ$ as a vector-valued mass flux.  One says that mass is conserved if the PDE
\begin{equation}
\frac{\partial\rho}{\partial t} + \Div \bJ = 0  \label{cont}
\end{equation}
applies.  Now suppose that $\Omega\subset\RR^3$ is any open and bounded set such that $\partial\Omega$ is continuously-differentiable.  Apply the divergence theorem to \eqref{cont} to show
\begin{equation}
\frac{d}{dt}\left(\int_\Omega \rho\,dm\right) = - \int_{\partial\Omega} \bJ\cdot\bn\,ds.  \label{conttwo}
\end{equation}
Interpret \eqref{conttwo} physically.  (\emph{What is the total mass inside $\Omega$?  How does that quantity change?  Relate to the net normal flux across the boundary.  Note that $\bn$ is the \emph{outward} normal direction.})
\end{document}
